\section{Bilinear sums and the sharp prime cutoff}
\label{sec:perron}

We now assemble the estimates, keeping track of the losses needed to pass
from smooth weights to the interval indicator. Throughout this section,
\(L=\log X\), \(AB\asymp X\),
\[
 H=\frac{B^2}{A},\qquad \mathcal Q=ABH^{1/3},
 \qquad A\mathcal Q=(AB)^{5/3}\asymp X^{5/3}.
\]
All support ratios are fixed. Coefficients are divisor-bounded, and their
fixed moment bounds may contribute fixed powers of \(L\). Write
\[
 \mathcal D(n)=G(n)-c_*\frac{\mu^2(n)}{N(n)^{1/6}}.
\]
Both terms vanish unless \(n\) is squarefree. Thus separate squarefree
restrictions may always be imposed on the two variables in a sum of
\(\mathcal D(ab)\); common factors are still permitted, since the complete
summand then vanishes. Let \(V\) denote a fixed smooth compactly supported
function on \((0,\infty)\).

The assembly uses the three ranges in Table~\ref{tab:height-ranges}.
The low range compares the Gauss sum with its model; the high ranges
treat the Gauss sum alone, after removing the undecomposed prime model.
The small constants and logarithmic powers are chosen below in an
explicit order.
\begin{table}[ht]
\centering
\small
\begin{tabular}{@{}lll@{}}
\toprule
Height range & Type I outer length & Type II shorter length\\
\midrule
$|t|\le L^{C_t}$ & $R\le X^{.40}$ & $cX^{1/3}\le B\ll X^{1/2}$\\[3pt]
$L^{C_t}\le T\ll X^{.01}$ & $R\le X^{.40}$ & $cX^{1/3}\le B\ll X^{1/2}$\\[3pt]
$X^{.01}\ll T\ll X^{1/6+\rho}$ & $R\le X^{1/3-\kappa/2}$
 & $X^{1/3-3\kappa}\le B\ll X^{1/2}$\\
\bottomrule
\end{tabular}
\caption{Ranges supplied by Proposition~\ref{prop:decomposition}, with
$AB\asymp X$ and fixed $c>0$. At high heights and $B\le X^{.40}$,
the cell parameter is $J=L^{D'}$ in the middle row and $J=X^\gamma$
in the last row. Larger $B$ uses averaged dispersion directly.}
\label{tab:height-ranges}
\end{table}

\subsection{Bilinear sums at low heights}

\begin{proposition}[Low-height bilinear comparison]
\label{prop:bilinear-low}
Suppose \(X^{.32}\le B\ll\sqrt X\), \(H\gg1\), and the coefficient
\(\beta_b\) satisfies the structural hypotheses of
Proposition~\ref{prop:dispersion}. At low heights assume, in addition,
its roughness and Siegel--Walfisz hypotheses. For every fixed \(D,U>0\),
the logarithmic parameters in that proposition can be chosen so that,
uniformly for \(|t|\le L^U\) and \(H\ge L^{C_g}\),
\begin{equation}
 \sum_{a,b}\alpha_a\beta_b\mathcal D(ab)N(ab)^{it}
                 V\!\left(\frac{N(ab)}X\right)
       \ll X^{5/6}L^{-D}.
 \label{eq:bilinear-low}
\end{equation}
There is also the following transition estimate. Suppose \(b\) is one
prime, \(a\) is a product of two primes, each of the three primes has
its own specified smooth dyad of length \(X^{1/3}L^{O(1)}\), and their
weights are bounded. If \(1\ll H<L^{C_g}\), then the left side of
\eqref{eq:bilinear-low} is
\begin{equation}
 O\!\left(X^{5/6}L^{-3/2}\right).
 \label{eq:bilinear-transition}
\end{equation}
The constants are uniform for the families of dyads in
Proposition~\ref{prop:decomposition}.
\end{proposition}

\begin{proof}
First omit the product envelope. In the notation of
Proposition~\ref{prop:dispersion}, multiplicativity gives
\[
 \sum_{a,b}\alpha_a\beta_bG(ab)N(ab)^{it}
   =\sum_a\alpha_aG(a)N(a)^{it}F_t(a).
\]
Apply Cauchy's inequality in \(a\) to the corrected estimate
\eqref{eq:dispersion-corrected}. Since
\(\sum_a|\alpha_a|^2\ll AL^{O(1)}\), its contribution is
\[
 \ll (A\mathcal Q)^{1/2}L^{-K/2+O(1)}
 \ll X^{5/6}L^{-K/2+O(1)}.
\]
The resulting main term is
\[
 c_*\sum_a\alpha_a\mu^2(a)N(a)^{it-1/6}
       \sum_b\beta_bN(b)^{it-1/6}.
\]
It remains to reinstate coprimality between \(a\) and \(b\).
Any common prime has norm greater than \(L^{C_r}\). The number of
pairs sharing such a prime is
\[
 \ll AB\sum_{N\mathfrak p>L^{C_r}}\frac1{N\mathfrak p^2}
 \ll AB L^{-C_r}.
\]
Here ideals of norm at most \(Y\) have count \(O(Y)\), so the last
inequality follows even on summing over all ideals. By Cauchy's
inequality and the coefficient moment bounds, the weighted number of
these pairs is \(\ll AB L^{-C_r/2+O(1)}\). Multiplication by
\(N(ab)^{-1/6}\) makes their model contribution
\(\ll X^{5/6}L^{-C_r/2+O(1)}\). Taking \(K,C_r\) large proves
\eqref{eq:bilinear-low} without its envelope.

For the transition, the specified prime dyads and their bounded weights
give
\[
 \sum_a|\alpha_a|^2\ll\frac A{L^2},\qquad
 \sum_b|\beta_b|^2\ll\frac BL,\qquad
 \sum_b|\beta_b|\ll\frac BL.
\]
The possible ordering multiplicity for the two primes in \(a\) is
bounded. Apply \eqref{eq:dispersion-short} and Cauchy's inequality.
The diagonal contributes
\[
 \left(\frac A{L^2}\frac{AB}L\right)^{1/2}
 \ll X^{5/6}H^{-1/6}L^{-3/2},
\]
and the remaining displayed main bound contributes
\[
 \left(\frac A{L^2}\frac{\mathcal Q}{L^2}\right)^{1/2}
 \ll X^{5/6}L^{-2}.
\]
The \(L^{-K}\) remainder is smaller by increasing \(K\).
Since \(H\gg1\), this bounds the Gauss-sum term as asserted.
The model term is bounded directly by
\(O(X^{5/6}L^{-3})\), using the independent counts on the three
specified prime dyads.

For completeness, the envelope introduces no logarithmic loss in these
arguments. Mellin inversion separates \(V(N(ab)/X)\) into common norm
twists with a rapidly decreasing, integrable weight. Keep shifts
\(|v|\le L^E\), for a sufficiently large fixed \(E\), in the
low-height estimates, choosing their permitted height exponent after
\(E\). From \(L^E\) up to \(X^c\), with a small fixed \(c>0\), the
coarse part of Proposition~\ref{prop:dispersion} and Cauchy's inequality
bound the uncorrected sum by \(X^{5/6}L^{O(1)}\); the model has the
same bound trivially. Rapid decay of the Mellin weight makes this tail
as small a log power as needed. Beyond \(X^c\), its rapid decay and
the trivial bound \(XL^{O(1)}\) give a power saving. This proves both
claims with the envelope.
\end{proof}

\subsection{Large heights and localization near the product boundary}

The next estimate treats the Gauss-sum term alone; the prime model will
be removed before decomposition. When \(B>X^{.40}\), the averaged
dispersion diagonal already saves a fixed power. For smaller \(B\),
divide both log-norm ranges into cells. Near the product boundary,
each cell has only boundedly many partners, so the localized diagonal
gains the square root of the cell width. Away from the boundary,
integration by parts in height supplies the saving.

Choose the absolute exponent in \eqref{eq:dispersion-power} so small
that \(0<\gamma<10^{-3}\), and then choose
\begin{equation}
 0<\kappa<\frac\gamma{100},\qquad
 0<\rho<\frac\kappa{100}.
 \label{eq:power-parameters}
\end{equation}
The decomposition uses a fixed \(0<\xi<\kappa/10\).
The availability of an absolute \(\gamma\) in the range
\(B\le X^{.40}\) is important: these choices precede, and do not
depend on, the number of prime factors permitted in the decomposition.

\begin{proposition}[Integrated bilinear estimate]
\label{prop:bilinear-height}
Let \(X_0\asymp X\), and let \(h_T\) be smooth, supported where
\(T\le |t|\le2T\), with
\(h_T^{(j)}(t)\ll_jT^{-j}\). Suppose
\[
 L^{C_t}\le T\ll X^{1/6+\rho},\qquad B\ll\sqrt X,
\]
and suppose that either
\begin{align*}
 T\ll X^{.01},&\qquad cX^{1/3}\le B,\quad c>0\text{ fixed},\tag{a}\\
 T\gg X^{.01},&\qquad X^{1/3-3\kappa}\le B.\tag{b}
\end{align*}
Assume the structural hypotheses of Proposition~\ref{prop:dispersion},
but impose neither roughness nor Siegel--Walfisz conditions.
For every fixed \(D>0\), by taking \(C_t\) sufficiently large,
\begin{equation}
 \sum_{a,b}\alpha_a\beta_bG(ab)V\!\left(\frac{N(ab)}X\right)
 \int h_T(t)\left(\frac{N(ab)}{X_0}\right)^{it}\frac{dt}{T}
       \ll X^{5/6}L^{-D}.
 \label{eq:bilinear-integrated}
\end{equation}
\end{proposition}

\begin{proof}
When \(B>X^{.40}\), apply Cauchy's inequality in \(a,t\), followed by
\eqref{eq:dispersion-height}. The diagonal costs
\[
 \ll A\sqrt B\,L^{O(1)}
 \asymp \frac X{\sqrt B}L^{O(1)}
 \ll X^{.80}L^{O(1)},
\]
whereas the error costs \(X^{5/6}L^{-K/2+O(1)}\).
The product envelope separates by Mellin inversion. Its shifts may
be truncated at a small fixed power of \(X\), since the tails are
negligible trivially. The retained shifts satisfy
\(2T+|u_0|\le X^{.17}\), because \(\rho<.001\).
This proves the assertion in this range.

Assume now \(B\le X^{.40}\). Partition the two fixed log-norm
intervals into disjoint cells of width \(1/J\), where
\[
 J=L^{D'}\quad\text{in case (a)},\qquad
 J=X^\gamma\quad\text{in case (b)}.
\]
There are \(O(J)\) cells on each side. Put
\[
 a_i=\|\alpha\,1_{\text{cell }i}\|_2,
 \qquad b_j=\|\beta\,1_{\text{cell }j}\|_2.
\]
A smooth majorant of an \(a\)-cell has derivatives of order \(s\)
bounded by \(O_s(J^s)\). Its lattice-point count is
\[
 O(A/J+\sqrt A+1)=O(A/J),
\]
because \(A\gg X^{.60}\) and \(\gamma<.30\).
Restricting \(\beta\) sharply to a cell preserves its divisor bounds
and squarefree support; there is no additional prime-polynomial
hypothesis in this range of \(B\).

For a fixed cell pair, Cauchy's inequality and the averaged dispersion
estimate give the bound
\begin{equation}
 C a_i\left(\frac A J b_j^2+\mathcal Q\varepsilon_J\right)^{1/2},
 \qquad
 \varepsilon_J=
 \begin{cases}
 L^{-K},&\text{in case (a)},\\
 X^{-10\gamma},&\text{in case (b)}.
 \end{cases}
 \label{eq:cell-pair}
\end{equation}
It holds uniformly after any fixed common Mellin shift: for
\(B\le X^{.40}\), Proposition~\ref{prop:dispersion} places no
restriction on \(u_0\). In particular it applies with the product
envelope, whose Mellin transform has bounded \(L^1\) norm.

Call a pair near if its rectangle meets
\(\left|\log(N(ab)/X_0)\right|\le c_0/J\), with a fixed
\(c_0>0\). Every cell has only \(O(1)\) near neighbors. Hence the
diagonal parts of \eqref{eq:cell-pair}, summed over near pairs, are
\begin{equation}
 \ll\sqrt{\frac A J}\,\|\alpha\|_2\|\beta\|_2
 \ll X^{5/6}L^{O(1)}H^{-1/6}J^{-1/2}.
 \label{eq:near-cell-diagonal}
\end{equation}
This is the required log saving in case (a), by choosing \(D'\)
large. In case (b), \(H\gg X^{-9\kappa}\), so its additional
power of \(X\) is at most
\(3\kappa/2-\gamma/2<0\).

The errors in \eqref{eq:cell-pair} are absolute errors; no reduction
with the mass of a \(b\)-cell is assumed. Even summing them over all
\(O(J^2)\) pairs gives at most
\begin{align}
 \sqrt{\mathcal Q\varepsilon_J}\,J\sum_i a_i
 &\ll \sqrt{\mathcal Q\varepsilon_J}\,J^{3/2}\|\alpha\|_2
       \notag\\
 &\ll X^{5/6}L^{O(1)}J^{3/2}\varepsilon_J^{1/2}.
 \label{eq:all-cell-error}
\end{align}
In case (a), take \(K\) large after \(D'\). In case (b), the
additional factor is \(X^{-7\gamma/2}\).

It remains to handle far pairs. Set
\(\ell=\log(N(ab)/X_0)\). On a far pair \(|\ell|>c_0/J\).
Integrating by parts \(q\) times in the height integral differentiates
only \(h_T\) and produces a factor \((J/T)^q\), a new height
function proportional to \(T^q h_T^{(q)}\), and the product weight
\((J\ell)^{-q}\). The new height function still has the same
normalized derivative bounds.

There is no loss by a power of \(J\) in the Mellin \(L^1\) norm.
Indeed, for a fixed
integer \(q\ge2\), choose a smooth function \(F_q\) on \(\mathbb R\)
which equals \(s^{-q}\) for \(|s|\ge c_0\) and is zero near zero.
Both \(F_q\) and \(F_q''\) are integrable. Consequently its Fourier
transform is integrable, and
\[
 \|\widehat{F_q(J\,\cdot)}\|_1=\|\widehat F_q\|_1<\infty.
\]
The smooth extension \(F_q(J\ell)\), multiplied by the fixed
envelope, therefore has Mellin \(L^1\) cost \(O_q(1)\), uniformly
in \(J\). All these shifts are independent of \(t\). Even a shift
larger than \(T\) is allowed here: it simply multiplies the fixed
Dirichlet-polynomial coefficients by phases, as reflected in the
unrestricted \(u_0\) assertion of Proposition~\ref{prop:dispersion}.

Apply \eqref{eq:cell-pair} again with this separated weight. Before
the integration-by-parts factor, the sum of all diagonal parts is
\begin{equation}
 \ll\sqrt{\frac A J}\left(\sum_i a_i\right)
                         \left(\sum_j b_j\right)
 \ll X^{5/6}L^{O(1)}H^{-1/6}J^{1/2}.
 \label{eq:far-cell-diagonal}
\end{equation}
In case (a), \((J/T)^q\le L^{-q(C_t-D')}\), which pays for this
loss by taking \(C_t\) sufficiently large. In case (b), it is at
most a constant times \(X^{-q(.01-\gamma)}\), which pays for
\(H^{-1/6}J^{1/2}\ll X^{3\kappa/2+\gamma/2}\).
The errors are already controlled by \eqref{eq:all-cell-error}.
This proves \eqref{eq:bilinear-integrated}.
\end{proof}

\subsection{Choice of the logarithmic parameters}

We record the order of choices to make the uniformity in the final
assembly explicit. First choose \(\gamma,\kappa,\rho\) as in
\eqref{eq:power-parameters}, then \(\xi\) and the fixed bin ratio of
Proposition~\ref{prop:decomposition}. Its number of terms is now at
most \(L^{C_{\mathrm{dec}}}\), and fix \(M\) larger than all the
coefficient moment exponents needed below. These exponents are
independent of \(C_r\).

Choose
\[
 D>C_{\mathrm{dec}}+10,\qquad
 D'>2(D+M+2),\qquad
 K>3D'+2(D+M+2).
\]
The log-width near-cell and error bounds are then respectively
\[
 X^{5/6}L^{M-D'/2}\qquad\text{and}\qquad
 X^{5/6}L^{M+3D'/2-K/2}.
\]
Choose \(C_t\) large enough for
the averaged dispersion estimate at derivative scale \(L^{D'}\)
and output power \(K\), for the far-cell estimate with a fixed
\(q\ge2\), and for the prime-model estimate below.

Only now choose the parameters for low heights. Take the required
corrected-variance output power larger than \(2(D+M+2)\), choose
the sieve depth and the logarithmic transform cutoffs in the proof
of Proposition~\ref{prop:dispersion}, and fix a height exponent
covering \(L^{C_t}\) together with all retained transform shifts.
Then choose \(C_r,C_g\) sufficiently large for the rough-support,
diagonal, and cube-approximation errors, and request the finitely
many Siegel--Walfisz saving powers needed for these choices.
That property is available for every fixed logarithmic conductor and
height range. Inserting \(L^{C_r}\) as a bin endpoint does not change
the exponent \(C_{\mathrm{dec}}\), and enlarging \(C_g\) only changes
fixed constants in the number of transition configurations.

Finally choose all arbitrary small-power losses below the positive
power gaps obtained with these fixed choices, including the sparse
support gap, which may depend on \(C_r\). There is no circular
dependence: the absolute power-width estimate is used only for
\(B\le X^{.40}\), where it is independent of prime-polynomial
complexity.

\subsection{A smooth truncation of the interval indicator}

\begin{lemma}[Fourier truncation]
\label{lem:perron-truncation}
Let \(1\le T\le X\), and let \(V\) be a fixed smooth function
supported in \((1/2,4)\), equal to one on a neighborhood of \([1,2]\).
There is a smooth function \(I_T\) such that
\begin{equation}
 I_T(s)=\frac1{2\pi}\int_{\mathbb R}
       \eta(t/T)\frac{1-2^{-it}}{it}e^{its}\,dt,
 \label{eq:interval-fourier}
\end{equation}
where \(\eta\) is a fixed smooth cutoff supported in \([-1,1]\)
and equal to one near zero. For every \(m>0\),
\begin{equation}
 \left|1_{[1,2]}(v)-V(v)I_T(\log v)\right|
 \ll_m \sum_{b\in\{0,\log2\}}
             (1+T|\log v-b|)^{-m}
 \label{eq:interval-boundary-error}
\end{equation}
on the support of \(V\), and the difference is zero outside it.
If coefficients collected at each integer norm satisfy
\(|c_n|\ll_\epsilon X^\epsilon\) for \(X/2<n<4X\), then
\begin{equation}
 \sum_n |c_n|\left|1_{[X,2X]}(n)
            -V(n/X)I_T(\log(n/X))\right|
       \ll_\epsilon X^\epsilon(1+X/T).
 \label{eq:perron-error}
\end{equation}
\end{lemma}

\begin{proof}
With the Fourier convention in \eqref{eq:interval-fourier}, the
Fourier transform of \(1_{[0,\log2]}\) is
\((1-2^{-it})/(it)\), interpreted continuously at zero.
The inverse transform of \(\eta(t/T)\) is a Schwartz approximate
identity \(T k(Ts)\) of integral one. Convolving the interval
indicator with it proves \eqref{eq:interval-fourier}.
Away from either endpoint, the error is bounded by the integral of
the rapidly decreasing kernel beyond that endpoint; this proves
\eqref{eq:interval-boundary-error}. In the portion where \(V\ne1\),
the distance to \([1,2]\) is bounded below, so the same estimate
holds there. On \(1/2\le v\le4\), distance in \(\log v\) is
comparable to distance in \(v\). Summing
\((1+T|n-X_0|/X)^{-2}\) over integers, for \(X_0=X,2X\),
costs \(O(1+X/T)\), uniformly in the location of \(X_0\).
This gives \eqref{eq:perron-error}, including exact endpoint equality.
\end{proof}

\subsection{Completion of the prime asymptotic}

\begin{proof}[Proof of Theorem~\ref{thm:main}]
Use Lemma~\ref{lem:perron-truncation} with
\[
 T_{\max}=X^{1/6+\rho}.
\]
At each norm the number of primary prime generators is divisor-bounded,
and \(|G(\pi)|=1\). Thus the lemma applies separately to the
Gauss-sum and model prime coefficients. Taking \(\epsilon<\rho\)
in \eqref{eq:perron-error}, their total truncation error is
\[
 \ll X^\epsilon(1+X/T_{\max})
   =o\!\left(\frac{X^{5/6}}L\right).
\]
It remains to estimate their difference with the smoothed indicator.

Split the integral in \eqref{eq:interval-fourier} into bounded
frequencies and smooth dyadic height blocks. Outside bounded
frequencies its two endpoint terms have the form
\[
 \int h_T(t)\left(\frac{N(n)}{X_0}\right)^{it}\frac{dt}T,
 \qquad X_0=X\text{ or }2X,
\]
where \(h_T^{(j)}\ll_jT^{-j}\). Indeed the factor \(1/(it)\)
is \(1/T\) times a smooth bounded function on a height block;
the cutoff at \(T_{\max}\) preserves these derivative bounds.
At bounded frequencies the original difference of endpoint terms
is bounded, including at zero.

For \(|t|\le L^{C_t}\), use the ordinary-block decomposition of
Proposition~\ref{prop:decomposition} on the difference
\(\mathcal D(n)\). The Type I terms are controlled by
Proposition~\ref{prop:typeI}; a norm twist of logarithmic height
only introduces logarithmic derivative bounds. The regular Type II
terms are controlled by Proposition~\ref{prop:bilinear-low} with
the arbitrary output power \(D\) fixed above. The sparse terms
have a fixed power saving.

The remaining terms are exactly the three-prime transition described
in Proposition~\ref{prop:decomposition}. They have bounded weights,
\(1\ll H<L^{C_g}\), and only
\(O((1+\log L)^{C})\) configurations for a fixed \(C\).
Proposition~\ref{prop:bilinear-low} bounds each by
\(X^{5/6}L^{-3/2}\). There are only \(O(1+\log L)\)
low-height blocks. Therefore their aggregate contribution is
\[
 \ll X^{5/6}L^{-3/2}(1+\log L)^{C+1}
   =o\!\left(\frac{X^{5/6}}L\right).
\]
All other low-height terms are smaller after summing their
\(L^{C_{\mathrm{dec}}}\) pieces.

At heights \(T\ge L^{C_t}\), first discard the \emph{undecomposed}
prime model. On the fixed envelope, write its norm polynomial as
\[
 P_0(t)=\sum_{n\asymp X}d_n n^{it},\qquad
 \sum_n|d_n|^2\ll X^{2/3}L^{O(1)}.
\]
The last bound follows by collecting the model coefficients
\(c_*N(\pi)^{-1/6}V(N(\pi)/X)\) at equal norms and using the
divisor moment bound. The ordinary mean-value inequality and
Cauchy's inequality give
\begin{align*}
 \int_{T\le|t|\le2T}|P_0(t)|\frac{dt}T
 &\ll X^{1/3}L^{O(1)}(1+X/T)^{1/2}\\
 &\ll X^{5/6}L^{O(1)}T^{-1/2},
\end{align*}
since \(T\ll X^{1/6+\rho}<X\). Taking \(C_t\) large makes
this negligible even after all \(O(L)\) height blocks. No
decomposition of this model term, and no cancellation between its
decomposed pieces, is required.

Apply the decomposition to the Gauss-sum term alone at the remaining
heights. Ordinary blocks have \(T\ll X^{.01}\); their no-stop
Type I terms satisfy \(R\le X^{.40}\). By
Proposition~\ref{prop:typeIheight}, the first term of their bound
has exponent at most
\[
 \frac12+\frac34(.40)+\frac12(.01)=.805<\frac56.
\]
For upper blocks \(T\gg X^{.01}\), use the upper-block
decomposition. Its no-stop terms have
\(R\le X^{1/3-\kappa/2}\), and their first-term exponent is
\[
 \frac12+\frac34\left(\frac13-\frac\kappa2\right)
       +\frac12\left(\frac16+\rho\right)
   =\frac56-\frac{3\kappa}{8}+\frac\rho2<\frac56.
\]
The pole term in Proposition~\ref{prop:typeIheight} is
\(X^{5/6}L^{O(1)}T^{-q}\) for every fixed \(q\), and is
therefore negligible above \(L^{C_t}\). These bounds permit the
uniformly smooth envelope at scale \(X/N(r)\), and any twist of
the outer coefficient has absolute value one.

Every nonsparse Type II term is in case (a) or case (b) of
Proposition~\ref{prop:bilinear-height}. Its bound can have the
arbitrary fixed output power \(D\). The sparse ordinary-block terms retain their
power saving uniformly in the norm twist. Summing the
\(O(L)\) height blocks and the at most
\(L^{C_{\mathrm{dec}}}\) pieces consequently costs less than
the chosen logarithmic saving. Changing the decomposition between
height blocks is legitimate because each is an exact identity on
squarefree arguments.

Combining the low and high heights with the truncation error proves
\begin{equation}
 \sum_{X<N(\pi)\le2X}
       \left(G(\pi)-c_*N(\pi)^{-1/6}\right)
       =o\!\left(\frac{X^{5/6}}{\log X}\right).
 \label{eq:prime-dyadic-comparison}
\end{equation}
Changing endpoint inclusion has only divisor-bounded cost.
Dyadically sum \eqref{eq:prime-dyadic-comparison}. For example,
on dyads above \(X^{1/2}\), the little-oh coefficient is uniformly
small, the logarithms are comparable to \(\log X\), and the
\(5/6\)-powers form a geometric sum. The terms below \(X^{1/2}\)
are \(O(X^{1/2})\) trivially. We obtain
\[
 \sum_{N(\pi)\le X}G(\pi)
   =c_*\sum_{N(\pi)\le X}N(\pi)^{-1/6}
       +o\!\left(\frac{X^{5/6}}{\log X}\right).
\]
Finally the prime ideal theorem in the fixed field gives
\(\#\{\pi:N(\pi)\le y\}\sim y/\log y\); omission of the
prime above \(3\) is immaterial. Partial summation yields
\[
 \sum_{N(\pi)\le X}N(\pi)^{-1/6}
   \sim\int_2^X\frac{u^{-1/6}}{\log u}\,du
   \sim\frac65\frac{X^{5/6}}{\log X}.
\]
This proves the asserted primary-prime asymptotic with constant
\(\frac65c_*\).
\end{proof}
